%=========
% Packages
%=========
 
\usepackage{
    amsmath,
    amsfonts,
    amssymb,
    amsthm,
    amscd,
    comment,
    enumitem,
    etoolbox,
    textcomp,% To avoid warning in gensymb
    gensymb,    % For \degree
    mathtools,
    mathdots,
    stmaryrd,
    graphicx% For \inplus
}
\usepackage[dvipsnames]{xcolor}
\usepackage{tikz-cd}
\usepackage{csquotes}
\usepackage{dsfont}
\usepackage{ytableau}
\usepackage{mathtools}

%================
% Layout Settings
%================

\usepackage[letterpaper,margin=1in]{geometry}
\setlength{\parindent}{0.5cm}

\clubpenalty=10000
\widowpenalty=10000
\displaywidowpenalty=10000

% \usepackage{enumitem}

% \newlist{enum}{enumerate}{1}
% \setlist[enum,1]{label=\arabic*), wide, labelwidth=!, labelindent=0pt, leftmargin=*}

\usepackage{pgfornament}
\usepackage[x11names]{xcolor}
\usepackage{pgfornament}
\usepackage{adjustbox}

\newcommand\orn[4]{%
\begin{center}
\fcolorbox{white}{#3}{{%
\color{#4}%
\pgfornament[width=15mm]{#1}
\pgfornament[width=15mm]{#2}
}}%
\end{center}}

%\orn{53}{54}{Plum}{White}


\usepackage{enumitem}

\newenvironment{enum}
  {\begin{enumerate}[label=\arabic*)]}
  {\end{enumerate}}

% wide, labelwidth=!, labelindent=0.5pt, leftmargin=*


%==============
% Font packages
%==============

\usepackage[T1]{fontenc}
\usepackage{lmodern}            % To get better quality text
\usepackage[colorlinks=true, linkcolor=blue, citecolor=blue, urlcolor=blue, breaklinks=true]{hyperref}

%============
% Referencing
%============

\usepackage{zref-clever}

\zcsetup{
  cap,                      % To capitalize names of theorems, etc.
  nameinlink=false,         % Options are true, false, single, and tsingle
  lastsep = {, and }        % Oxford comma
}

\zcRefTypeSetup{assumption}{
    Name-sg = Assumption ,
    name-sg = assumption ,
    Name-pl = Assumptions ,
    name-pl = assumptions ,
}
\zcRefTypeSetup{cor}{
    Name-sg = Corollary ,
    name-sg = corollary ,
    Name-pl = Corollaries ,
    name-pl = corollaries ,
}
\zcRefTypeSetup{defin}{
    Name-sg = Definition ,
    name-sg = definition ,
    Name-pl = Definitions ,
    name-pl = definitions ,
}
\zcRefTypeSetup{enumi}{
    Name-sg = {} ,
    name-sg = {} ,
    Name-pl = {} ,
    name-pl = {} ,
}
\zcRefTypeSetup{equation}{
    Name-sg = {} ,
    name-sg = {} ,
    Name-pl = {} ,
    name-pl = {} ,
}
\zcRefTypeSetup{eg}{
    Name-sg = Example ,
    name-sg = example ,
    Name-pl = Examples ,
    name-pl = examples ,
}
\zcRefTypeSetup{lem}{
    Name-sg = Lemma ,
    name-sg = lemma ,
    Name-pl = Lemmas ,
    name-pl = lemmas ,
}
\zcRefTypeSetup{prop}{
    Name-sg = Proposition ,
    name-sg = proposition ,
    Name-pl = Propositions ,
    name-pl = propositions ,
}
\zcRefTypeSetup{rem}{
    Name-sg = Remark ,
    name-sg = remark ,
    Name-pl = Remarks ,
    name-pl = remarks ,
}
\zcRefTypeSetup{theo}{
    Name-sg = Theorem ,
    name-sg = theorem ,
    Name-pl = Theorems ,
    name-pl = theorems ,
}

% Map counters → type “item”
\zcsetup{
  countertype = {
    enumi=item, enumii=item, enumiii=item, enumiv=item
  },
% Tell zref-clever how nested counters reset
  counterresetby = {
    enumii=enumi, enumiii=enumii, enumiv=enumiii
  }
}
\zcRefTypeSetup{item}{
    Name-sg = {} ,
    name-sg = {} ,
    Name-pl = {} ,
    name-pl = {} ,
}

% Map the *counter* "theo" to different *types* as each environment starts:
\AddToHook{env/assumption/begin}{\zcsetup{countertype={theo=assumption}}}
\AddToHook{env/cor/begin}{\zcsetup{countertype={theo=cor}}}
\AddToHook{env/defin/begin}{\zcsetup{countertype={theo=defin}}}
\AddToHook{env/eg/begin}{\zcsetup{countertype={theo=eg}}}
\AddToHook{env/lem/begin}{\zcsetup{countertype={theo=lem}}}
\AddToHook{env/prop/begin}{\zcsetup{countertype={theo=prop}}}
\AddToHook{env/rem/begin}{\zcsetup{countertype={theo=rem}}}

%---------------------------------------------
% To enable compatibility with cleveref syntax
%---------------------------------------------

\newcommand{\cref}[1]{\zcref{#1}}
\newcommand{\Cref}[1]{\zcref[S]{#1}}
\newcommand{\cpageref}[1]{\zcpageref{#1}}

%----------------------------------
% To add links to description lists
% Syntax:
% \begin{description}
%     \item[\namedlabel{R1}{R1}] Item
% \end{description}
% Refer to \ref{R1}.
%----------------------------------

\makeatletter
\def\namedlabel#1#2{\begingroup
    (#2)%
    \def\@currentlabel{(#2)}%
    \phantomsection\label{#1}\endgroup
}
\makeatother

\renewcommand{\descriptionlabel}[1]{%
  \hspace\labelsep \upshape\bfseries #1%
}

%=================
% Math Definitions
%=================

\newcommand\C{\mathbb{C}}
\newcommand\N{\mathbb{N}}
\newcommand\Q{\mathbb{Q}}
\newcommand\R{\mathbb{R}}
\newcommand\Z{\mathbb{Z}}
\newcommand{\Sc}{\mathbb{S}}
\newcommand\kk{\Bbbk}

\newcommand{\gS}{\mathfrak{S}}

%---------------
% Colors
%---------------
\definecolor{FieldsGreen}{HTML}{2b7443}

%---------------
% Math Operators
%---------------

\DeclareMathOperator{\Aut}{Aut}
\DeclareMathOperator{\End}{End}
\DeclareMathOperator{\Hom}{Hom}
\DeclareMathOperator{\id}{id}
\DeclareMathOperator{\im}{im}                   % image of a map
\DeclareMathOperator{\Mat}{Mat}

%---------------
% Math Objects
%---------------
\newcommand{\supp}{\mathrm{supp}}
\newcommand{\Par}{\mathrm{Par}}
\newcommand{\SPar}{\mathrm{SPar}}
\newcommand{\WComp}{\mathrm{WComp}}
\newcommand{\SComp}{\mathrm{SComp}}
\newcommand{\Pol}{\mathrm{Pol}}
\newcommand{\LPol}{\mathrm{LPol}}
\newcommand{\Sym}{\mathrm{Sym}}
\newcommand{\LSym}{\mathrm{LSym}}
\newcommand{\LAlt}{\mathrm{LAlt}}
\newcommand{\Alt}{\mathrm{Alt}}
\newcommand{\Si}{\mathrm{Sign}}
\newcommand{\sgn}{\mathrm{sgn}}
\newcommand{\GL}{\mathrm{GL}}

%=====
% TikZ
%=====

\usepackage{tikz}
\usetikzlibrary{arrows.meta,patterns}
\usetikzlibrary{shapes.geometric,decorations.markings,decorations.pathreplacing}
\usepackage{tikz-cd}

%-------
% Styles
%-------

\tikzset{% Automatically centers a TikZ diagram vertically
    anchorbase/.style={
        >=To,
        line cap = round, line join = round,
        baseline={([yshift=-0.5ex]current bounding box.center)},
    }
}
\tikzset{% Syntax: \begin{tikzpicture}[centerzero={0,0.2}]
    centerzero/.style={
        >=To,
        line cap = round, line join = round,
        baseline={([yshift=-0.5ex]#1)},
        },
    centerzero/.default={0,0}
}

%=====================
% Theorem Environments
%=====================

\newtheorem{theo}{Theorem}[section]
\newtheorem{prop}[theo]{Proposition}
\newtheorem{lem}[theo]{Lemma}
\newtheorem{cor}[theo]{Corollary}
\newtheorem{conj}[theo]{Conjecture}

\theoremstyle{definition}
\newtheorem{assumption}[theo]{Assumption}
\newtheorem{defin}[theo]{Definition}
\newtheorem{rem}[theo]{Remark}
\newtheorem{eg}[theo]{Example}

\numberwithin{equation}{section}
\allowdisplaybreaks

\renewcommand{\theenumi}{\alph{enumi}}
\setenumerate[1]{label=(\alph*)}          % Only need to reference enumerate items with \ref{}

\setcounter{tocdepth}{2}

%========
% Toggles
%========

\newtoggle{comments}
\newtoggle{details}
\newtoggle{detailsnote}

% \toggletrue{comments}   % To include comments (default is false)
% \toggletrue{details}   % To include details (default is false)
% \toggletrue{detailsnote}   % For including note about details toggle (default is false)

\iftoggle{comments}{%
    \usepackage[notref,notcite]{showkeys}   % To show labels
    \newcommand{\acomments}[1]{% For comments from Alistair
        \ \\
        {\color{red}
            \textbf{AS:} #1
        }
        \ \\
    }
    \newcommand{\ycomments}[1]{% For comments from Yaolong
        \ \\
        {\color{red}
            \textbf{YS:} #1
        }
        \ \\
    }
    \newcommand{\hcomments}[1]{% For comments from Hugo
        \ \\
        {\color{cyan}
            \textbf{HF:} #1
        }
        \ \\
    }
    \newcommand{\ecomments}[1]{% For comments from Elena
        \ \\
        {\color{cyan}
            \textbf{EK:} #1
        }
        \ \\
    }
    \newcommand{\tcomments}[1]{% For comments from Teresa
        \ \\
        {\color{cyan}
            \textbf{TL:} #1
        }
        \ \\
    }
    \newcommand{\mcomments}[1]{% For comments from Mario
        \ \\
        {\color{cyan}
            \textbf{MTD:} #1
        }
        \ \\
    }
    }{%
        \newcommand{\acomments}[1]{}
        \newcommand{\ycomments}[1]{}
        \newcommand{\hcomments}[1]{}
        \newcommand{\ecomments}[1]{}
        \newcommand{\tcomments}[1]{}
        \newcommand{\mcomments}[1]{}
    }

\definecolor{neongreen}{RGB}{57, 255, 20}

\iftoggle{details}{%
    \newcommand{\details}[1]{
        \ \\
        {\color{OliveGreen}
            \textbf{Details:} #1
        }
        \\
    }
}{%
    \newcommand{\details}[1]{\ignorespaces}
} 


\usepackage{wasysym}
\renewcommand\qedsymbol{
\fbox{}
} 

%\DeclareEmphSequence{\color{blue}\itshape}

    \newcommand{\change}[1]{
        {\color{Plum}
            #1
        }
    }
\newcommand{\plumstrip}{%
  \begin{tikzpicture}
    \fill[violet] (0,0) rectangle (\linewidth,0.3);
  \end{tikzpicture}%
}

\newcommand{\orangestrip}{%
  \begin{tikzpicture}
    \fill[orange] (0,0) rectangle (\linewidth,0.3);
  \end{tikzpicture}%
}

\newcommand{\notdone}[1]{
        {\color{Orange}
            #1
        }
    }

\newcommand{\ceqref}[1]{% For putting an equation number over an equals sign, \cref{<label>}
    \mathrel{\mathop{=}\limits^{\text{\cref{#1}}}}
}

\newcommand{\ceqsref}[2]{% For putting equation numbers over and below an equals sign
    \mathrel{\mathop{=}\limits^{\text{\cref{#1}}}_{\text{\cref{#2}}}}
}

\makeatletter
%%%%%%%%%%%%%%%%%%%%%%%%%%%%%%%%%%%%%%%%%%%%%%%%%%%%%%%%%%%%%%%%%%%%%%%%%
%% Equals signs carrying equation numbers, for use in aligned displays.
%%
%%   \eqalignref [<extra>]{<top label>}
%%   \eqalignsref[<extra>]{<top label>}{<bottom label>}
%%
%% IMPORTANT: these macros CONTAIN the alignment tab. Write
%%
%%     a &= b \\
%%     c \eqalignref{eq:foo} d \\
%%     e \eqalignsref{eq:foo}{eq:bar} f
%%
%% and NOT "c &\eqalignref{eq:foo} d".
%%
%% The label overhangs the "=" by half its excess width on each side.
%% That overhang is measured and compensated automatically, so (5.14)
%% and (4.3) receive the same visual gap. The optional <extra>
%% (default 0mu) is a minimum additional separation added on each
%% side, on top of the ordinary \thickmuskip that any relation gets;
%% [0mu] reproduces exactly the spacing of a plain "&=".
%%%%%%%%%%%%%%%%%%%%%%%%%%%%%%%%%%%%%%%%%%%%%%%%%%%%%%%%%%%%%%%%%%%%%%%%%

\newlength{\eqalignref@lab}
\newlength{\eqalignref@tmp}

% Measure one label at script size; keep the running maximum.
% Global, because each alignment cell is a group: a length set before
% the "&" would otherwise be restored before the "&" is reached.
\newcommand{\eqalignref@measure}[1]{%
    \settowidth{\eqalignref@tmp}{%
        \ensuremath{\scriptstyle\text{\cref{#1}}}%
    }%
    \ifdim\eqalignref@tmp>\eqalignref@lab
        \global\eqalignref@lab=\eqalignref@tmp
    \fi
}

% Turn the max label width into the padding needed on each side.
\newcommand{\eqalignref@finish}{%
    \settowidth{\eqalignref@tmp}{\ensuremath{=}}%
    \global\advance\eqalignref@lab -\eqalignref@tmp
    \ifdim\eqalignref@lab<\z@
        \global\eqalignref@lab=\z@
    \fi
    \global\eqalignref@lab=.5\eqalignref@lab
}

\newcommand{\eqalignref@pad}[1]{%
    \mspace{#1}\hspace{\eqalignref@lab}%
}

\newcommand{\eqalignref}[2][0mu]{%
    \global\eqalignref@lab=\z@
    \eqalignref@measure{#2}%
    \eqalignref@finish
    \eqalignref@pad{#1}%
    &\mathrel{\mathop{=}\limits^{%
        \mathclap{\text{\cref{#2}}}%
    }}%
    \eqalignref@pad{#1}%
}

\newcommand{\eqalignsref}[3][0mu]{%
    \global\eqalignref@lab=\z@
    \eqalignref@measure{#2}%
    \eqalignref@measure{#3}%
    \eqalignref@finish
    \eqalignref@pad{#1}%
    &\mathrel{\mathop{=}\limits^{%
        \mathclap{\text{\cref{#2}}}%
    }_{%
        \mathclap{\text{\cref{#3}}}%
    }}%
    \eqalignref@pad{#1}%
}
\makeatother
